\section{Controller Rules and Execution Semantics}
\label{app:rules}
This appendix describes the public v0.0.40 source anchor, commit
\path{37487f35610076c1016e1b59d3bf982388d1a275}.
The run bundle does not include a source commit for each job. Consequently, this release anchor documents the distributed implementation, but is not a per-job build attestation.

\subsection{Modality Confidence and Route Selection}
The scorer is implemented in \path{libs/ktem/ktem/reasoning/mara_route_scorer.py}; cost adjustments and tie breaking are in \path{mara_route_costing.py} in the same directory. For modality $m$, its probe confidence is zero if no evidence is present. Otherwise it is the clipped and four-decimal-rounded value
\begin{align}
 c_m=\operatorname{clip}_{[0,1]}\big(&b_m+0.04\min(n_m,3)\nonumber\\
 &+0.14s_m+0.08d_m+0.16\ell_m\nonumber\\
 &+i_m+t_m-0.25(1-h_m)\big),
\end{align}
where $n_m$ is evidence count, $s_m$ is the bounded top score, $d_m$ is its margin over the second score, $\ell_m$ is source/page locator coverage, and $h_m$ is a backend-health indicator. These features are proxies, not calibrated probabilities.

The bases $(b_T,b_V,b_E,b_G)$ are $(0.45,0.35,0.35,0.30)$. Visual, element, and graph intent add $0.12$, $0.15$, and $0.20$ to the corresponding modality. Text/OCR availability adds $0.08$ for text or $0.05$ for visual/element evidence; graph has no such bonus. If every probe is empty, the implementation inserts a text prior with one item, locator quality $0.7$, and text availability; visual intent can also insert a visual prior. These defaults must not be read as observed retrieval success.

Text/OCR intent terms, task type, and scope drive heuristic features. Examples include visual terms such as \emph{figure}, \emph{chart}, and \emph{layout}; element terms such as \emph{table}, \emph{row}, and \emph{cell}; and calculation terms such as \emph{ratio}, \emph{average}, and \emph{sum}. Graph intent includes summary/compare/study-guide tasks or corresponding question terms.

Element coverage is eligible when evidence is nonempty, locator quality is at least $0.5$, and text/OCR is present. Hybrid eligibility requires text confidence at least $0.45$, either visual or element confidence at least $0.45$, and no graph intent. On MMDocRAG without visual intent and with text confidence at least $0.65$, effective visual confidence is capped at $0.45$.

Using these effective confidences, the initial quality proxies are
\begin{align*}
 \widehat Q_T&=0.15+c_T,\\
 \widehat Q_V&=0.10+c_V,\\
 \widehat Q_E&=0.10+c_E,\\
 \widehat Q_G&=0.05+c_G,\\
 \widehat Q_H&=0.20+(c_T+\max(c_V,c_E))/2.
\end{align*}
The element and hybrid expressions start at zero when their eligibility conditions fail. Table~\ref{tab:quality-rules} lists subsequent additions. Quality values are then floored at zero, with four-decimal rounding.

\begin{table*}[t]
\centering
\small
\begin{tabularx}{\textwidth}{@{}Xrrrrr@{}}
\toprule
Condition & Text & Page & Element & Graph & Hybrid \\
\midrule
Visual intent & $-0.08$ & $+0.32$ & 0 & 0 & 0 \\
Element intent & 0 & 0 & $+0.25$ if eligible & 0 & $+0.10$ if eligible \\
Structured calculation & $+0.10$ & 0 & $+0.15$ if eligible & 0 & $+0.12$ if eligible \\
Graph intent present / absent & 0 & 0 & 0 & $+0.35$ / $-0.25$ & 0 \\
$c_T\geq0.65$ and no visual intent & $+0.25$ & $-0.20$ & 0 & 0 & $-0.10$ \\
MMDocRAG and no visual intent & $+0.20$ & $-0.20$ & 0 & 0 & $-0.12$ \\
Finance and structured calculation & $+0.20$ & 0 & 0 & 0 & $-0.05$ \\
\bottomrule
\end{tabularx}
\caption{Additive quality-proxy rules. Conditions can accumulate. Dataset-family rules are explicit implementation choices, not learned generalization.}
\label{tab:quality-rules}
\end{table*}

The base cost penalties are reported in Section~\ref{sec:controller}. Visual intent reduces page-image cost by $0.18$; eligible element intent reduces element cost by $0.05$. MMDocRAG adds $0.22$ to page-image and $0.18$ to hybrid cost. A missing visual generator adds $0.05$ to page-image cost with visual intent, or $0.25$ without it, and adds $0.12$ to hybrid cost. Costs are floored at zero and rounded to four decimals.

The scorer skips ineligible element retrieval. It skips page-image and hybrid routes on MMDocRAG when $c_T\geq0.65$ and visual intent is absent. It also skips either of those routes when cost is at least $0.55$ and its quality proxy does not exceed the text proxy. Among allowed, non-skipped candidates it maximizes $\max(0,\widehat Q-\widehat C)$. Equal scores prefer text, page-image, element, hybrid, then global graph. If no candidate remains, it falls back to permitted text or the first permitted route. The automatic benchmark route list contains those five evidence routes, rather than the separate direct-answer and guarded comparison configurations.

\subsection{Adequacy, Verification, and Recovery}
The relevant sources are \path{libs/ktem/ktem/docqa/controller.py}, \path{retrieval_adequacy.py}, \path{verification.py}, and \path{execution.py} in that DocQA directory, together with \path{libs/ktem/ktem/reasoning/mara.py}.

For ordinary retrieval, captured concrete evidence yields a \emph{good} decision unless a domain adequacy rule fails. Metadata without concrete evidence is \emph{ambiguous}; absent evidence is \emph{poor}. Finance checks look for required statement fields and, depending on origin, page-scoped support. These are permissive heuristics, not a universal numeric retrieval-quality threshold. For example, a finance request can pass when some required field aliases are present even if others are absent; benchmark-origin finance requests relax the page-scoping requirement.

Verification first checks domain support, contradictions, semantic rules, and token overlap. The general lexical path accepts overlap of at least $\min(2,\text{number of claim tokens})$, with additional short-evidence and source-summary rules. For claims that do not pass support checks, the unsupported-confidence threshold is $0.85$ in light mode, $0.75$ in strict mode, and $0.90$ in the finance domain. These thresholds act on heuristic values: explicit domain rejection gives $1.0$, contradiction $0.95$, certain unsupported summaries $0.85$, a strict directional-conflict condition $0.78$, and missing evidence tokens $0.80$; the remaining default is $0.55$. None is an independently calibrated probability. Unsupported claims request revision; insufficient evidence can request retry or abstention, and guardrail handling determines the returned answer.

There are two retrieval mechanisms. The pipeline-level \path{_should_retry_retrieval} returns true only for thorough mode with no retrieved documents. In the structured controller, \path{_retrieve_and_evaluate} repeats an ambiguous retrieval once. If retrieval remains non-good, \path{_switch_after_failed_retrieval} iterates alternative permitted routes until one yields good evidence. Each candidate can itself undergo the ambiguous-retrieval repeat. The result records the successful alternative; if none succeeds, execution returns the guarded result for the failed initial decision. Therefore the switch flag counts recorded successful recovery, not all attempted routes. The release does not implement the blanket assertion that total execution is capped at one retrieval retry.

\subsection{Visual Evidence Versus Visual Generation}
The source \path{libs/ktem/ktem/reasoning/mara_visual_answering.py} can return an OCR-extracted short answer before calling a VLM. For benchmark-origin requests, it searches visual evidence text for short answer-bearing patterns. With no visual generator it can return an evidence-only message or let another generation path handle the request, depending on configuration. The synthesis column \path{generator_backend} records configured identities such as \path{Qwen/Qwen3-8B} and \path{local_qwen3_vl}; the metric CSV does not contain the per-answer \path{generation_backend} or \path{visual_generation_gate} fields needed to identify these paths. A future generator comparison must export those fields and generated answers before attributing differences to a VLM.

\section{Coverage and Complete Diagnostics}
\label{app:results}
\begin{table}[t]
\centering
\small
\begin{tabular}{@{}lrrr@{}}
\toprule
Dataset & Questions & Sets & Records \\
\midrule
QASPER & 200 & 4 & 600 \\
ALCE-ASQA & 200 & 4 & 600 \\
MMDocRAG & 120 & 20 & 600 \\
FinanceBench & 150 & 3 & 600 \\
RAGTruth & 300 & 6 & 900 \\
SlideVQA & 120 & 4 & 240 \\
\midrule
Total & 1,090 & 41 & 3,540 \\
\bottomrule
\end{tabular}
\caption{An artifact set is a completed run/shard directory with prediction and report outputs. One set can contain multiple route configurations. Questions count unique dataset--example identifiers, and records count configuration-specific outputs.}
\label{tab:coverage}
\end{table}

\subsection{Metrics and Interpretation}
Tables~\ref{tab:all-quality} and \ref{tab:all-operations} include every dataset/configuration row in the released aggregate table. Missing metrics are shown as dashes, not zeros. Citation metadata recall and inline recall are kept separate. Metadata can include evidence attached by an adapter and must not be interpreted as verified support for every generated claim. Unsupported-claim rates come from the system's own verifier. False abstention is the harness flag against a nonempty reference answer, not a human assessment of answerability.

F1 and EM score the harness's normalized \path{answer_for_scoring}, which can differ from the user-facing answer. The synthesis does not include either answer string. SlideVQA's two configurations tie on all six published paired quality/locator metrics, despite 25 distinct per-example F1 values. Their mean recorded generation times are 4.61 and 2.55 seconds. This limits generator comparison even though routing probes are observable. The RAGTruth rows likewise belong to the local document-QA adaptation; token F1 is not the official hallucination-detection evaluation.

The instrumented \path{total_seconds} values have separate index and generation fields, while parse/retrieval fields in this bundle are zero. Zero instrumented retrieval time must not be interpreted as free retrieval. These logs do not establish full interactive startup latency, comparable cache state, or monetary/energy cost. The table's p95 is the sorted value at rounded index $0.95(n-1)$, as in the released synthesis script. Runtime scoring penalties in Section~\ref{sec:controller} are unrelated units.

\begin{table*}[t]
\centering\small
\begin{tabular}{@{}llrrrrrrrr@{}}
\toprule
Dataset & Config. & $n$ & F1 & EM & $C_{meta}$ & $C_{inline}$ & Page & FA & Native \\
\midrule
QASPER & Auto & 200 & 21.43 & 0.50 & 89.65 & 89.65 & -- & 0.00 & 21.43 \\
QASPER & Guarded & 200 & 20.06 & 0.50 & 89.65 & 89.65 & -- & 6.00 & 20.06 \\
QASPER & Text & 200 & 21.44 & 0.50 & 90.91 & 90.91 & -- & 1.50 & 21.44 \\
\midrule
ALCE-ASQA & Auto & 200 & 25.74 & 13.00 & 95.50 & 95.50 & -- & 0.00 & 77.07 \\
ALCE-ASQA & Guarded & 200 & 15.16 & 7.50 & 95.50 & 95.50 & -- & 46.00 & 63.20 \\
ALCE-ASQA & Text & 200 & 26.68 & 14.50 & 95.50 & 95.50 & -- & 0.00 & 77.14 \\
\midrule
MMDocRAG & Auto & 120 & 34.60 & 0.83 & 67.99 & 25.12 & 82.50 & 0.00 & 34.60 \\
MMDocRAG & Element & 120 & 5.57 & 0.00 & 3.26 & 2.85 & 5.83 & 76.67 & 5.57 \\
MMDocRAG & Hybrid & 120 & 33.99 & 0.00 & 30.77 & 13.44 & 58.33 & 0.00 & 33.99 \\
MMDocRAG & Page+VLM & 120 & 16.38 & 0.00 & 48.41 & 25.53 & 56.67 & 0.00 & 16.38 \\
MMDocRAG & Text & 120 & 35.44 & 0.00 & 68.62 & 24.03 & 81.67 & 5.83 & 35.44 \\
\midrule
FinanceBench & Auto & 150 & 15.10 & 0.67 & 51.00 & 18.22 & 60.00 & 2.67 & 8.67 \\
FinanceBench & Guarded & 150 & 15.33 & 0.67 & 51.00 & 18.22 & 60.00 & 2.67 & 8.67 \\
FinanceBench & Direct & 150 & 17.34 & 1.33 & 0.00 & 0.00 & 0.00 & 0.00 & 2.67 \\
FinanceBench & Text & 150 & 14.83 & 0.67 & 48.56 & 13.67 & 60.00 & 7.33 & 8.67 \\
\midrule
RAGTruth & Auto & 300 & 21.68 & 0.00 & 55.67 & 99.00 & -- & 0.00 & 54.00 \\
RAGTruth & Guarded & 300 & 21.45 & 0.00 & 55.67 & 99.00 & -- & 1.00 & 54.00 \\
RAGTruth & Text & 300 & 23.06 & 0.00 & 98.33 & 98.33 & -- & 0.00 & 54.00 \\
\midrule
SlideVQA & Auto & 120 & 43.25 & 0.00 & 100.00 & 78.75 & 100.00 & 0.00 & 43.25 \\
SlideVQA & Page+VLM & 120 & 43.25 & 0.00 & 100.00 & 78.75 & 100.00 & 0.00 & 43.25 \\
\bottomrule
\end{tabular}
\caption{All released quality/locator rows, on a 0--100 scale. $C_{meta}$/$C_{inline}$ are metadata/inline citation-locator recall; FA is false abstention. Native is the local dataset-adapted score, not an official external evaluation. Page+VLM denotes a configured visual generator, not proof of its use on every answer. SlideVQA rows are retained as non-discriminating diagnostics.}
\label{tab:all-quality}
\end{table*}

\begin{table*}[t]
\centering\small
\begin{tabular}{@{}llrrrrrrr@{}}
\toprule
Dataset & Config. & $n$ & UCR & Switch & Errors & Timeout & Mean (s) & p95 (s) \\
\midrule
QASPER & Auto & 200 & 0.00 & 0.00 & 0 & 0 & 1.40 & 2.49 \\
QASPER & Guarded & 200 & 6.00 & 0.00 & 0 & 0 & 1.44 & 2.55 \\
QASPER & Text & 200 & 0.00 & 0.00 & 0 & 0 & 1.70 & 2.92 \\
\midrule
ALCE-ASQA & Auto & 200 & 0.00 & 0.00 & 0 & 0 & 1.90 & 2.81 \\
ALCE-ASQA & Guarded & 200 & 46.00 & 0.00 & 0 & 0 & 2.00 & 2.96 \\
ALCE-ASQA & Text & 200 & 0.00 & 0.00 & 0 & 0 & 2.05 & 3.16 \\
\midrule
MMDocRAG & Auto & 120 & 0.00 & 5.00 & 0 & 0 & 22.37 & 101.54 \\
MMDocRAG & Element & 120 & 0.00 & 0.00 & 0 & 0 & 10.05 & 34.07 \\
MMDocRAG & Hybrid & 120 & 0.00 & 0.00 & 0 & 0 & 32.34 & 76.62 \\
MMDocRAG & Page+VLM & 120 & 0.00 & 0.00 & 0 & 0 & 29.56 & 70.63 \\
MMDocRAG & Text & 120 & 0.00 & 0.00 & 0 & 0 & 14.33 & 38.48 \\
\midrule
FinanceBench & Auto & 150 & 15.33 & 2.67 & 0 & 0 & 3.59 & 11.50 \\
FinanceBench & Guarded & 150 & 14.00 & 2.67 & 0 & 0 & 3.48 & 11.59 \\
FinanceBench & Direct & 150 & 0.00 & 0.00 & 0 & 0 & 1.17 & 0.87 \\
FinanceBench & Text & 150 & 14.67 & 0.00 & 0 & 0 & 15.41 & 29.05 \\
\midrule
RAGTruth & Auto & 300 & 0.00 & 0.00 & 3 & 0 & 0.88 & 3.16 \\
RAGTruth & Guarded & 300 & 1.00 & 0.00 & 3 & 0 & 0.88 & 3.18 \\
RAGTruth & Text & 300 & 0.00 & 0.00 & 5 & 0 & 2.94 & 5.06 \\
\midrule
SlideVQA & Auto & 120 & 0.00 & 0.00 & 0 & 0 & 4.61 & 5.01 \\
SlideVQA & Page+VLM & 120 & 0.00 & 0.00 & 0 & 0 & 2.55 & 2.86 \\
\bottomrule
\end{tabular}
\caption{All released operational rows. UCR is heuristic unsupported-claim rate and Switch is the recorded route-switch rate (both percentages). Error and timeout columns are counts. Mean/p95 are instrumented per-record seconds under the retained runs and their differing budgets; they are not matched-budget interactive latency measurements.}
\label{tab:all-operations}
\end{table*}

\begin{table*}[t]
\centering\small
\begin{tabular}{@{}llrrlrrr@{}}
\toprule
Dataset & Comparator & $n$ & Auto minus comparator & 95\% CI & W & L & T \\
\midrule
QASPER & Guarded & 200 & 1.38 & [0.28, 2.68] & 8 & 2 & 190 \\
QASPER & Text & 200 & -0.01 & [-0.40, 0.39] & 6 & 9 & 185 \\
ALCE-ASQA & Guarded & 200 & 10.57 & [7.30, 13.98] & 55 & 3 & 142 \\
ALCE-ASQA & Text & 200 & -0.94 & [-2.22, 0.08] & 1 & 5 & 194 \\
MMDocRAG & Element & 120 & 29.03 & [25.11, 32.89] & 102 & 4 & 14 \\
MMDocRAG & Hybrid & 120 & 0.61 & [-3.79, 4.97] & 45 & 49 & 26 \\
MMDocRAG & Page+VLM & 120 & 18.22 & [13.19, 22.96] & 82 & 18 & 20 \\
MMDocRAG & Text & 120 & -0.84 & [-4.02, 2.32] & 45 & 54 & 21 \\
FinanceBench & Guarded & 150 & -0.23 & [-1.03, 0.41] & 2 & 2 & 146 \\
FinanceBench & Direct & 150 & -2.24 & [-7.54, 2.89] & 58 & 48 & 44 \\
FinanceBench & Text & 150 & 0.27 & [-2.20, 2.64] & 23 & 29 & 98 \\
RAGTruth & Guarded & 300 & 0.22 & [0.00, 0.50] & 3 & 0 & 297 \\
RAGTruth & Text & 300 & -1.38 & [-2.41, -0.30] & 90 & 129 & 81 \\
SlideVQA & Page+VLM & 120 & 0.00 & [0.00, 0.00] & 0 & 0 & 120 \\
\bottomrule
\end{tabular}
\caption{All 14 released paired comparisons, with F1 differences on a 0--100 scale. Intervals use 1,000 resamples of existing paired records, with no multiple-comparison adjustment. Positive intervals for MMDocRAG Element and Page+VLM comparisons do not isolate routing from generation or configuration effects.}
\label{tab:all-ci}
\end{table*}


\subsection{Oracle Definitions and Pairing}
Let $a_i$ be controller F1 and $b_i$ the largest observed F1 among fixed comparison configurations for question $i$. The bundle excludes \path{controller_auto} when constructing $b_i$. We distinguish the signed difference of means
\[
 D=\frac1n\sum_i(b_i-a_i)
\]
from its recorded positive-part regret
\[
 R_+=\frac1n\sum_i\max(b_i-a_i,0).
\]
The two differ when auto beats every fixed comparison on some examples. MMDocRAG has $D=8.66$ and $R_+=11.05$ on a 0--100 F1 scale. Neither is a clean estimate of recoverable routing loss because fixed configurations can differ in generators, guards, and execution budgets. On FinanceBench the best comparison is direct-answer for 79 questions and guarded for 22; neither is an automatic evidence-route option. We therefore report this as a diagnostic envelope and avoid a claimed reachable-oracle percentage.

Bootstrap analysis joins by dataset and example ID. Each resample draws $n$ paired differences with replacement; 1,000 resampled means are sorted and the rounded-index 2.5th and 97.5th percentiles are used. The seed is reset to 20260705 for each comparison, matching the release script. Wins and losses are based on rounded four-decimal paired differences. These intervals describe the fixed existing sample and do not measure robustness across seeds, model families, or independent reruns.

\subsection{Execution Provenance and Resource Limits}
The input manifest identifies two A100 GPUs per shard for QASPER, ALCE-ASQA, FinanceBench, and RAGTruth, and two L40S GPUs for MMDocRAG and SlideVQA. The common text job timeout field is 180 seconds. MMDocRAG text and element jobs record 360 seconds, hybrid and page-image jobs 600 seconds, and controller jobs 360 seconds; SlideVQA records 360 seconds. These manifest timeout fields must not be conflated with the route-level timeout, which defaults to 90 seconds for auto/guarded in the distributed route configuration.

All four MMDocRAG controller shards use a distinct manifest named
\path{mmdocrag-dev15-stat120-controller-t360.routes.json};
other MMDocRAG configurations use
\path{mmdocrag-dev15-stat120.routes.json}.
The controller run names identify an extended-timeout repair. The synthesis lists manifest paths but does not embed their effective JSON, so their exact setting differences cannot be independently reconstructed from this ZIP. One text shard also comes from a missing-PDF repair run. Zero recorded route timeouts in the retained runs therefore does not demonstrate equivalent timeout tolerance.

The synthesis report is dated 5 July 2026 and is distributed as a v0.0.40 release asset. It includes aggregate metrics, paired deltas, bootstrap intervals, oracle rows, win/loss/tie counts, per-example metric records, the run ledger, input manifest, and synthesis script. Original raw predictions, effective route manifests, and per-job source/build identities are not included. The released tables can be reproduced from their metric records; exact end-to-end reruns require additional source data and configuration provenance. Separate post-hoc RAGTruth prompt-budget diagnostics are not pooled into these numbers.

\section{Artifact Access and Demonstration Deployment}
\label{app:artifact}
The public release is
\url{https://github.com/262412/MARA/releases/tag/v0.0.40}.
Its assets include the updated Windows installer \path{slide-app.zip}, its SHA-256 file,
\path{mara-full-system-benchmark-synthesis-v0.0.40.zip},
and the historical release's source archive. The repository is
\url{https://github.com/262412/MARA};
the video is
\url{https://youtu.be/owRaHCzSVNg}.
These are also the links to enter in the submission form.

\subsection{Reviewer Installation and Source Identity}
For Windows x64, download \path{slide-app.zip} from the release page, extract it to a writable directory, and run \path{Install.cmd} followed by \path{Start.cmd}. The installer provisions Python 3.10.19 and pinned dependencies in an independent environment; it requires internet access but no existing Python, Git, or compiler for the included UI/CLI runtime. Configure both a chat model and an embedding model in the Web UI's resources tab before indexing the included sample document. No model weights or credentials are supplied.

This replacement installer, published on 23 September 2026, is built from main plus installation fixes at commit \texttt{e6afa5dc}; its packages report 0.0.41 and its manifest records the full source commit and file hashes. It is not a rebuild of the historical benchmark executable. Isolated Windows installation, CLI checks, HTTP/UI startup, repeat installation preserving configuration, and checksum-failure handling were tested. A new real-model index/query run and a cross-platform installation matrix have not been established.

The v0.0.40 tag remains the historical source anchor, at commit \texttt{37487f35}. For source inspection, it can still be retrieved with the following shell commands:
\begin{verbatim}
git clone --branch v0.0.40 \
  https://github.com/262412/MARA.git
cd MARA
\end{verbatim}
Replacing a release asset does not change this tag. A fresh Windows check of the historical \texttt{uv sync --extra mara} command failed while building \path{llama-cpp-python==0.2.7}, pulled in by \path{kotaemon[all]}, because \texttt{nmake} and C/C++ compilers were unavailable. It is therefore not the reviewer quick-start route. The ZIP runtime omits that optional in-process backend; other local or visual routes require their own documented backend configuration.

The historical README's \texttt{pip install mara-research-cli} route was unavailable on PyPI and TestPyPI when checked on 22 September 2026. The reviewer ZIP installs its included application wheels directly. Apache-2.0 licensing and Kotaemon attribution are retained; the historical NOTICE still uses the earlier product name \emph{Slides}. Minimum model-serving RAM/VRAM requirements have not been measured.

After ZIP installation and provider configuration, representative Windows commands are:
\begin{verbatim}
MARA.cmd app doctor --json
MARA.cmd docqa index ./example.pdf
MARA.cmd docqa files
MARA.cmd docqa ask --help
\end{verbatim}
The package README supplies an example question and its source document. Provider health checks do not replace a successful end-to-end index/query check. Users should inspect citations and status, and use an appropriate source document with permission to process it. Local-only and external-API configurations have different hardware and data-flow requirements.

\subsection{Video and Benchmark Identity}
The unchanged video illustrates source selection, page preview, text questions, citations, and the reasoning-status card. Its user deployment uses Deepseek/Google APIs and a packaged UI label of 0.0.18. The benchmark identifies Qwen3-8B/bge-m3 and the repository release v0.0.40. These are deployment and distribution identifiers, not interchangeable claims about one executable. In the released code, the displayed app version may come from \path{KH_APP_VERSION} or installed package metadata; a UI version string alone is not a source-commit identifier. We do not infer binary equivalence between the recording and benchmark release.

The recording is a user-workflow demonstration, not a visual-generator comparison or an offline-privacy demonstration. The paper's multimodal and recovery limitations remain applicable even when the interface reports a successful text answer.
